\begin{lemma}[Second derivative along a circle]
Let $n\in\mathbb{N}$, let $P\in\mathbb{R}[X_1,\dots,X_n]$, let $\mathbb{R}^n$
carry the Euclidean inner product, and let $u:\mathbb{R}^n\to\mathbb{R}$ be the
polynomial function $u(y)=P(y_1,\dots,y_n)$. Let $C,A,B\in\mathbb{R}^n$ and set
\[
  \gamma(s)=C+\cos(s)\,A+\sin(s)\,B\qquad(s\in\mathbb{R}).
\]
Then for every $t\in\mathbb{R}$ the function $s\mapsto u(\gamma(s))$ has second
derivative at $t$ given by
\[
  \frac{d^2}{ds^2}\Big|_{s=t}u(\gamma(s))
  =D^2u\bigl(\gamma(t)\bigr)\bigl[\gamma'(t),\gamma'(t)\bigr]
   -\bigl\langle\nabla u(\gamma(t)),\,\cos(t)A+\sin(t)B\bigr\rangle ,
\]
where $\gamma(t)=C+\cos(t)A+\sin(t)B$ and $\gamma'(t)=-\sin(t)A+\cos(t)B$.
\end{lemma}
\begin{proof}
Write $E$ for $\mathbb{R}^n$ with the Euclidean inner product, and
$f(s)=u(\gamma(s))=P(\gamma(s))$ for $s\in\mathbb{R}$. Coordinatewise,
\begin{equation}\label{eq:gcsd-coord}
  \gamma(s)_i=C_i+\cos(s)A_i+\sin(s)B_i\qquad(1\le i\le n).
\end{equation}

\emph{Step 1: $\gamma$ is twice differentiable, with
$\gamma''=-(\gamma-C)$.} By \eqref{eq:gcsd-coord} each coordinate function
$s\mapsto\gamma(s)_i$ is a constant plus a linear combination of $\cos$ and
$\sin$, hence infinitely differentiable, with
\[
  \frac{d}{ds}\gamma(s)_i=-\sin(s)A_i+\cos(s)B_i,
  \qquad
  \frac{d^2}{ds^2}\gamma(s)_i=-\cos(s)A_i-\sin(s)B_i,
\]
using $\cos'=-\sin$, $\sin'=\cos$ and that the constant $C_i$ differentiates to
$0$. A map into the finite-dimensional space $E$ is differentiable at a point
exactly when each of its $n$ coordinate functions is, and its derivative is then
assembled coordinatewise; applying this twice gives that $\gamma$ is twice
differentiable with
\begin{equation}\label{eq:gcsd-gamma}
  \gamma'(s)=-\sin(s)A+\cos(s)B,
  \qquad
  \gamma''(s)=-\bigl(\cos(s)A+\sin(s)B\bigr).
\end{equation}

\emph{Step 2: $u$ is differentiable along $\gamma$, with
$f'(s)=\sum_{i=1}^{n}(\partial_iP)(\gamma(s))\,\gamma'(s)_i$.} Fix
$s\in\mathbb{R}$. By \noderef{MvPolyFDeriv} applied to $P$ at the point
$\gamma(s)$, the polynomial function $v\mapsto P(v)$ on the coordinate space
$\mathbb{R}^n$ is Frechet differentiable at $\gamma(s)$ with derivative
$w\mapsto\sum_{i=1}^{n}(\partial_iP)(\gamma(s))w_i$. The identity map of
$\mathbb{R}^n$, read from the Euclidean norm to the coordinate norm, is a
continuous linear isomorphism, because any two norms on a finite-dimensional
real vector space are equivalent; composing with it, the chain rule gives that
$u$ is Frechet differentiable at $\gamma(s)$ with
\[
  Du(\gamma(s))(w)=\sum_{i=1}^{n}(\partial_iP)(\gamma(s))\,w_i\qquad(w\in E).
\]
Since $\gamma$ is differentiable at $s$ by Step 1, the chain rule applied to
$f=u\circ\gamma$ gives that $f$ is differentiable at $s$ with
\begin{equation}\label{eq:gcsd-first}
  f'(s)=Du(\gamma(s))\bigl(\gamma'(s)\bigr)
      =\sum_{i=1}^{n}(\partial_iP)(\gamma(s))\,\gamma'(s)_i .
\end{equation}

\emph{Step 3: each coefficient in \eqref{eq:gcsd-first} is differentiable in
$s$.} Fix $1\le i\le n$ and $s\in\mathbb{R}$. Repeating Step 2 verbatim with $P$
replaced by the polynomial $\partial_iP$ --- that is, by \noderef{MvPolyFDeriv}
applied to $\partial_iP$ at $\gamma(s)$, the same norm-equivalence transfer, and
the chain rule along $\gamma$ --- the function $s\mapsto(\partial_iP)(\gamma(s))$
is differentiable with
\begin{equation}\label{eq:gcsd-coeff}
  \frac{d}{ds}(\partial_iP)(\gamma(s))
  =\sum_{j=1}^{n}\bigl(\partial_j\partial_iP\bigr)(\gamma(s))\,\gamma'(s)_j .
\end{equation}

\emph{Step 4: differentiating \eqref{eq:gcsd-first}.} The right-hand side of
\eqref{eq:gcsd-first} is a sum of $n$ products, the $i$-th being
$(\partial_iP)(\gamma(s))\cdot\gamma'(s)_i$. Its first factor is differentiable
in $s$ by Step 3 and its second factor is differentiable in $s$ by Step 1, with
derivative $\gamma''(s)_i$. By the product rule and then the sum rule, $f'$ is
differentiable and
\begin{align*}
  f''(s)&=\sum_{i=1}^{n}\left(
      \Bigl(\frac{d}{ds}(\partial_iP)(\gamma(s))\Bigr)\gamma'(s)_i
      +(\partial_iP)(\gamma(s))\,\gamma''(s)_i\right)\\
   &=\sum_{i=1}^{n}\sum_{j=1}^{n}
       \bigl(\partial_j\partial_iP\bigr)(\gamma(s))\,\gamma'(s)_j\,\gamma'(s)_i
     +\sum_{i=1}^{n}(\partial_iP)(\gamma(s))\,\gamma''(s)_i ,
\end{align*}
using \eqref{eq:gcsd-coeff} and splitting the resulting finite sum into two.

\emph{Step 5: the double sum is the Hessian term.} By \noderef{MvPolyHessian}
applied at the point $\gamma(s)$ with $v=w=\gamma'(s)$,
\[
  D^2u(\gamma(s))\bigl[\gamma'(s),\gamma'(s)\bigr]
  =\sum_{i=1}^{n}\sum_{j=1}^{n}\gamma'(s)_i\,\gamma'(s)_j\,
     \bigl(\partial_i\partial_jP\bigr)(\gamma(s)).
\]
Interchanging the names of the two summation indices in the double sum of
Step 4 --- a relabelling of a finite double sum --- turns it into exactly this
expression.

\emph{Step 6: the single sum is the gradient term.} By
\noderef{MvPolyGradient} applied at the point $\gamma(s)$, the $i$-th coordinate
of $\nabla u(\gamma(s))$ is $(\partial_iP)(\gamma(s))$ for every $i$. Hence, by
the coordinate formula for the Euclidean inner product,
\[
  \sum_{i=1}^{n}(\partial_iP)(\gamma(s))\,\gamma''(s)_i
  =\sum_{i=1}^{n}\bigl(\nabla u(\gamma(s))\bigr)_i\,\gamma''(s)_i
  =\bigl\langle\nabla u(\gamma(s)),\,\gamma''(s)\bigr\rangle .
\]
By \eqref{eq:gcsd-gamma} we have $\gamma''(s)=-\bigl(\cos(s)A+\sin(s)B\bigr)$,
so, the inner product being linear in its second argument,
\[
  \sum_{i=1}^{n}(\partial_iP)(\gamma(s))\,\gamma''(s)_i
  =-\bigl\langle\nabla u(\gamma(s)),\,\cos(s)A+\sin(s)B\bigr\rangle .
\]

\emph{Step 7: conclusion.} Substituting Steps 5 and 6 into Step 4 and
specialising to $s=t$,
\[
  f''(t)=D^2u\bigl(\gamma(t)\bigr)\bigl[\gamma'(t),\gamma'(t)\bigr]
   -\bigl\langle\nabla u(\gamma(t)),\,\cos(t)A+\sin(t)B\bigr\rangle .
\]
By Steps 2 and 4 the function $f$ is twice differentiable, so its second
derivative at $t$ is $f''(t)$; and $\gamma(t)=C+\cos(t)A+\sin(t)B$ together with
$\gamma'(t)=-\sin(t)A+\cos(t)B$ by \eqref{eq:gcsd-gamma}. This is the assertion.
\end{proof}
